% =============================================================================
% BÁO CÁO KẾT QUẢ THỰC NGHIỆM
%
% HƯỚNG DẪN UPLOAD LÊN OVERLEAF:
%   1. Tạo thư mục "figures/" trong project Overleaf
%   2. Copy 12 file PNG vào thư mục figures/ với tên như bên dưới:
%
%   results/antwerp/plots/vae_loss.png       → figures/vae_loss_antwerp.png
%   results/lorawan/plots/vae_loss.png       → figures/vae_loss_lorawan.png
%   results/indoor/plots/vae_loss.png        → figures/vae_loss_indoor.png
%
%   results/antwerp/plots/dt_estimation.png  → figures/dt_estimation_antwerp.png
%   results/lorawan/plots/dt_estimation.png  → figures/dt_estimation_lorawan.png
%   results/indoor/plots/dt_estimation.png   → figures/dt_estimation_indoor.png
%
%   results/antwerp/plots/roc_comparison.png → figures/roc_comparison_antwerp.png
%   results/lorawan/plots/roc_comparison.png → figures/roc_comparison_lorawan.png
%   results/indoor/plots/roc_comparison.png  → figures/roc_comparison_indoor.png
%
%   results/antwerp/plots/madrl_rewards.png  → figures/madrl_rewards_antwerp.png
%   results/lorawan/plots/madrl_rewards.png  → figures/madrl_rewards_lorawan.png
%   results/indoor/plots/madrl_rewards.png   → figures/madrl_rewards_indoor.png
% =============================================================================

\documentclass[12pt,a4paper]{article}

\usepackage[utf8]{inputenc}
\usepackage[T5]{fontenc}
\usepackage[vietnamese]{babel}
\usepackage{geometry}
\geometry{margin=2.5cm}

\usepackage{graphicx}
\usepackage{booktabs}
\usepackage{array}
\usepackage{multirow}
\usepackage{amsmath}
\usepackage{amssymb}
\usepackage{hyperref}
\usepackage{xcolor}
\usepackage{caption}
\usepackage{subcaption}
\usepackage{float}
\usepackage{siunitx}

\hypersetup{
    colorlinks=true,
    linkcolor=blue,
    citecolor=blue,
    urlcolor=blue,
}

\graphicspath{{figures/}}

\title{\textbf{Kết Quả Thực Nghiệm}\\
\large Khung Đề Xuất: VAE + Digital Twin (EKF) + MADRL+GAT\\
Đánh Giá Trên 3 Bộ Dữ Liệu IoT}
\author{}
\date{}

\begin{document}
\maketitle
\tableofcontents
\newpage

% =============================================================================
\section{Thiết Lập Thực Nghiệm}
% =============================================================================

\subsection{Bộ Dữ Liệu}

Thực nghiệm được tiến hành trên ba bộ dữ liệu IoT thực tế có đặc trưng khác nhau
về quy mô, công nghệ truyền thông và môi trường triển khai:

\begin{itemize}
    \item \textbf{Sigfox Antwerp}: Dữ liệu mạng Sigfox thu thập tại thành phố Antwerp,
    Bỉ. Bộ dữ liệu gồm 14.377 mẫu với không gian đặc trưng VAE 84 chiều (84 trạm
    gốc). Nhãn chất lượng gốc mất cân bằng nghiêm trọng (chỉ 1,9\% mẫu tốt),
    nên ngưỡng thích nghi được áp dụng tại $-123{,}9$ dBm (trung vị RSSI),
    đưa tỷ lệ về $50{,}0\%$.

    \item \textbf{LoRaWAN Italy}: Dữ liệu mạng LoRaWAN thu thập tại Ý, quy mô lớn
    với 190.502 mẫu và 5 đặc trưng. Ngưỡng phân loại tại trung vị
    $-69{,}0$ dBm, cho tỷ lệ $47{,}9\%$ mẫu chất lượng tốt.

    \item \textbf{Indoor WiFi}: Dữ liệu định vị WiFi trong nhà với 1.100 mẫu
    và 5 đặc trưng. Ngưỡng phân loại tại $-63{,}9$ dBm,
    tỷ lệ mẫu chất lượng tốt là $49{,}5\%$.
\end{itemize}

\subsection{Phương Pháp Tiêm Lỗi (Fault Injection)}

Để tạo nhãn chuẩn (ground truth) cho bài toán phát hiện bất thường, 5\% mẫu
trong mỗi chuỗi quan sát được tiêm lỗi nhân tạo theo ba kiểu. Biên độ tiêm
được giới hạn ở $\pm 2\sigma$ và $+1{,}5\sigma$ — tương ứng với các sự kiện
thực tế trong mạng IoT (nhiễu giao thoa nhẹ, suy hao tín hiệu, cảm biến
trôi giá trị), thay vì các cực trị thống kê quá rõ ràng:

\begin{itemize}
    \item \textbf{Spike} ($+2\sigma$): RSSI tăng vừa phải, mô phỏng nhiễu giao thoa.
    \item \textbf{Drop} ($-2\sigma$, số trạm gốc về 0): Suy hao tín hiệu và mất kết nối.
    \item \textbf{Stuck} ($+1{,}5\sigma$, số trạm gốc về 0): Cảm biến trôi giá trị nhẹ.
\end{itemize}

Số lỗi tiêm vào lần lượt là: Antwerp 718/14.377, LoRaWAN 9.525/190.502,
Indoor 55/1.100.

\subsection{Phương Pháp So Sánh (Baseline)}

\begin{itemize}
    \item \textbf{K-Means} ($k=2$): Phân cụm không giám sát cho phân loại chất lượng.
    \item \textbf{GMM} ($k=2$): Mô hình hỗn hợp Gaussian cho phân loại chất lượng.
    \item \textbf{LOF} ($n\_neighbors=20$): Local Outlier Factor cho phát hiện bất thường.
    \item \textbf{Rolling Z-Score} (cửa sổ 20 bước): Phát hiện bất thường dựa trên
    thống kê cục bộ trên chuỗi RSSI.
\end{itemize}

Lưu ý: K-Means và GMM chỉ thực hiện phân loại chất lượng; LOF và Rolling
Z-Score chỉ thực hiện phát hiện bất thường. Khung đề xuất DT-EKF thực hiện
\textit{đồng thời cả hai nhiệm vụ}.

% =============================================================================
\section{Kết Quả và Phân Tích}
% =============================================================================

\subsection{Nén Ngữ Nghĩa bằng VAE}

VAE được huấn luyện 30 epoch để học biểu diễn tiềm ẩn gọn gàng từ không gian
đặc trưng thô. Bảng~\ref{tab:vae} tổng hợp kết quả:

\begin{table}[H]
\centering
\caption{Kết quả huấn luyện VAE trên ba bộ dữ liệu}
\label{tab:vae}
\begin{tabular}{lcccccc}
\toprule
\textbf{Bộ dữ liệu} & \textbf{Chiều vào} & \textbf{Chiều tiềm ẩn}
    & \textbf{Tỷ lệ nén} & \textbf{MSE tái tạo} & \textbf{MAE} & \textbf{KL} \\
\midrule
Antwerp  & 84 & 16 & $5{,}25\times$ & 0,008213 & 0,0531 & 0,0012 \\
LoRaWAN  & 5  & 3  & $1{,}67\times$ & 0,015992 & 0,0891 & 0,0000 \\
Indoor   & 5  & 3  & $1{,}67\times$ & 0,053747 & 0,1734 & 0,0001 \\
\bottomrule
\end{tabular}
\end{table}

% -----------------------------------------------------------------------
% ẢNH 1: VAE Training Loss
% File cần có trong figures/:
%   vae_loss_antwerp.png, vae_loss_lorawan.png, vae_loss_indoor.png
% -----------------------------------------------------------------------
\begin{figure}[H]
    \centering
    \begin{subfigure}[b]{0.32\textwidth}
        \includegraphics[width=\textwidth]{vae_loss_antwerp}
        \caption{Antwerp (84$\to$16)}
    \end{subfigure}
    \hfill
    \begin{subfigure}[b]{0.32\textwidth}
        \includegraphics[width=\textwidth]{vae_loss_lorawan}
        \caption{LoRaWAN (5$\to$3)}
    \end{subfigure}
    \hfill
    \begin{subfigure}[b]{0.32\textwidth}
        \includegraphics[width=\textwidth]{vae_loss_indoor}
        \caption{Indoor WiFi (5$\to$3)}
    \end{subfigure}
    \caption{Đường cong loss huấn luyện VAE. Đường xanh: loss tái tạo;
    đường đỏ: KL divergence. Cả hai thành phần hội tụ nhanh sau vài epoch đầu.}
    \label{fig:vae_loss}
\end{figure}

Trên Antwerp (84 chiều vào), VAE đạt MSE tái tạo = 0,0082, cho thấy biểu diễn
tiềm ẩn 16 chiều bảo toàn tốt thông tin tín hiệu đa trạm. Trên LoRaWAN và
Indoor (5 chiều vào), tỷ lệ nén $1{,}67\times$ vẫn mang lại không gian tiềm
ẩn có cấu trúc hơn, loại bỏ nhiễu đo lường và cung cấp đầu vào chất lượng
cao cho module MADRL.

% -----------------------------------------------------------------------

\subsection{Ước Lượng Trạng Thái và Phân Loại của Digital Twin}

\subsubsection{Ước Lượng Trạng Thái bằng EKF}

Digital Twin sử dụng EKF để ước lượng liên tục trạng thái kênh truyền
gồm RSSI trung bình và số trạm gốc hoạt động. Hình~\ref{fig:dt_estimation}
so sánh giá trị quan sát thực tế (màu xám) với ước lượng EKF (màu màu)
trên 500 bước đầu tiên.

% -----------------------------------------------------------------------
% ẢNH 2: Digital Twin State Estimation
% File cần có: dt_estimation_antwerp.png, dt_estimation_lorawan.png, dt_estimation_indoor.png
% -----------------------------------------------------------------------
\begin{figure}[H]
    \centering
    \begin{subfigure}[b]{\textwidth}
        \includegraphics[width=\textwidth]{dt_estimation_antwerp}
        \caption{Antwerp (Sigfox) — RSSI-RMSE = 4,834 dBm; BS-RMSE = 7,022}
    \end{subfigure}
    \vspace{0.3cm}
    \begin{subfigure}[b]{\textwidth}
        \includegraphics[width=\textwidth]{dt_estimation_lorawan}
        \caption{LoRaWAN (Italy) — RSSI-RMSE = 7,045 dBm; BS-RMSE = 0,118}
    \end{subfigure}
    \vspace{0.3cm}
    \begin{subfigure}[b]{\textwidth}
        \includegraphics[width=\textwidth]{dt_estimation_indoor}
        \caption{Indoor WiFi — RSSI-RMSE = 2,443 dBm; BS-RMSE = 1,234}
    \end{subfigure}
    \caption{Ước lượng trạng thái Digital Twin (EKF). Đường xám: tín hiệu
    quan sát (có chứa lỗi tiêm); đường màu: ước lượng EKF. EKF làm mịn
    nhiễu và phục hồi xu hướng sau mỗi lỗi tiêm.}
    \label{fig:dt_estimation}
\end{figure}

RSSI-RMSE LoRaWAN cao nhất (7,045 dBm) do dải dao động tín hiệu LoRaWAN rộng
hơn ($-20$ đến $-120$ dBm) so với Sigfox và WiFi — xấp xỉ tuyến tính hóa
cục bộ của EKF gặp hạn chế trong môi trường tín hiệu phi tuyến mạnh. Ngược
lại, Indoor đạt RSSI-RMSE thấp nhất (2,443 dBm) nhờ tín hiệu WiFi ổn định
hơn trong không gian hẹp.

\subsubsection{Phân Loại Chất Lượng và Phát Hiện Bất Thường}

Bảng~\ref{tab:dt_metrics} tổng hợp kết quả DT-EKF cho hai tác vụ. Ngưỡng
phân loại được xác định qua tiêu chí Youden's J trên RSSI ước lượng, cho
phép DT-EKF thích nghi với đặc trưng từng bộ dữ liệu mà không cần hiệu chỉnh
thủ công.

\begin{table}[H]
\centering
\caption{Kết quả DT-EKF: phân loại chất lượng và phát hiện bất thường}
\label{tab:dt_metrics}
\begin{tabular}{lcccc}
\toprule
\multirow{2}{*}{\textbf{Bộ dữ liệu}}
    & \multicolumn{2}{c}{\textbf{Phân loại chất lượng}}
    & \multicolumn{2}{c}{\textbf{Phát hiện bất thường}} \\
\cmidrule(lr){2-3} \cmidrule(lr){4-5}
 & \textbf{Accuracy} & \textbf{AUC} & \textbf{RSSI-RMSE (dBm)} & \textbf{AUC} \\
\midrule
Antwerp  & 0,7403 & 0,7948 & 4,834 & 0,8653 \\
LoRaWAN  & 0,7825 & 0,8614 & 7,045 & 0,8337 \\
Indoor   & 0,9345 & 0,9725 & 2,443 & 0,9784 \\
\bottomrule
\end{tabular}
\end{table}

% -----------------------------------------------------------------------
% ẢNH 3: ROC Curves
% File cần có: roc_comparison_antwerp.png, roc_comparison_lorawan.png, roc_comparison_indoor.png
% -----------------------------------------------------------------------
\begin{figure}[H]
    \centering
    \begin{subfigure}[b]{\textwidth}
        \includegraphics[width=\textwidth]{roc_comparison_antwerp}
        \caption{Antwerp --- Chất lượng: DT-EKF (0,795) \ensuremath{>} GMM (0,697)
        \ensuremath{>} K-Means (0,676);
        Bất thường: Z-Score (0,920) \ensuremath{>} DT-EKF (0,865)
        \ensuremath{\gg} LOF (0,506)}
    \end{subfigure}
    \vspace{0.3cm}
    \begin{subfigure}[b]{\textwidth}
        \includegraphics[width=\textwidth]{roc_comparison_lorawan}
        \caption{LoRaWAN --- Chất lượng: DT-EKF (0,861) \ensuremath{\gg}
        GMM (0,516) \ensuremath{\approx} K-Means (0,505);
        Bất thường: Z-Score (0,919) \ensuremath{>} DT-EKF (0,834)
        \ensuremath{\gg} LOF (0,500)}
    \end{subfigure}
    \vspace{0.3cm}
    \begin{subfigure}[b]{\textwidth}
        \includegraphics[width=\textwidth]{roc_comparison_indoor}
        \caption{Indoor --- Chất lượng: DT-EKF (0,973) \ensuremath{>}
        K-Means (0,830) \ensuremath{>} GMM (0,666);
        Bất thường: DT-EKF (0,978) \ensuremath{>} Z-Score (0,958)
        \ensuremath{\gg} LOF (0,489)}
    \end{subfigure}
    \caption{Đường cong ROC so sánh các phương pháp. \textit{Trái}: phân loại
    chất lượng (K-Means nét đứt xanh dương; GMM nét đứt xanh lá; DT-EKF nét
    liền đỏ). \textit{Phải}: phát hiện bất thường (LOF nét đứt xanh dương;
    Z-Score nét đứt xanh lá; DT-EKF nét liền đỏ).}
    \label{fig:roc}
\end{figure}

% -----------------------------------------------------------------------

\subsection{Phân Bổ Tài Nguyên Hợp Tác bằng MADRL+GAT}

Module MADRL+GAT được huấn luyện trong 50 episode với 5 tác nhân tương tác
qua đồ thị chú ý (GAT). Không gian trạng thái gọn (3 chiều với LoRaWAN và
Indoor, 16 chiều với Antwerp) cho phép các tác nhân học chính sách ổn định
trong số episode hạn chế. Hình~\ref{fig:madrl} thể hiện quá trình học.

% -----------------------------------------------------------------------
% ẢNH 4: MADRL Rewards
% File cần có: madrl_rewards_antwerp.png, madrl_rewards_lorawan.png, madrl_rewards_indoor.png
% -----------------------------------------------------------------------
\begin{figure}[H]
    \centering
    \begin{subfigure}[b]{0.32\textwidth}
        \includegraphics[width=\textwidth]{madrl_rewards_antwerp}
        \caption{Antwerp\\Eval: 1916,01}
    \end{subfigure}
    \hfill
    \begin{subfigure}[b]{0.32\textwidth}
        \includegraphics[width=\textwidth]{madrl_rewards_lorawan}
        \caption{LoRaWAN\\Eval: 2519,26}
    \end{subfigure}
    \hfill
    \begin{subfigure}[b]{0.32\textwidth}
        \includegraphics[width=\textwidth]{madrl_rewards_indoor}
        \caption{Indoor WiFi\\Eval: 3088,95}
    \end{subfigure}
    \caption{Đường cong phần thưởng tích lũy MADRL+GAT (đường mờ: thô;
    đường đậm: làm mịn). Phần thưởng tăng đơn điệu trên cả ba bộ dữ liệu.}
    \label{fig:madrl}
\end{figure}

Phần thưởng đánh giá (evaluation reward) sau 20 episode kiểm tra tăng theo
quy mô và đặc tính dữ liệu: 1.916 (Antwerp, 14K mẫu,
tín hiệu Sigfox đa chiều) $\to$ 2.519 (LoRaWAN, 190K mẫu) $\to$ 3.089
(Indoor, tín hiệu WiFi ổn định, biên độ dao động nhỏ hơn).

% -----------------------------------------------------------------------

\subsection{So Sánh Tổng Thể}

Bảng~\ref{tab:summary} tổng hợp toàn bộ kết quả so sánh.

\begin{table}[H]
\centering
\caption{Bảng tổng hợp kết quả thực nghiệm trên ba bộ dữ liệu IoT}
\label{tab:summary}
\resizebox{\textwidth}{!}{%
\begin{tabular}{llccccccc}
\toprule
\multirow{2}{*}{\textbf{Bộ dữ liệu}} &
\multirow{2}{*}{\textbf{Phương pháp}} &
\multicolumn{2}{c}{\textbf{Phân loại chất lượng}} &
\multicolumn{3}{c}{\textbf{Phát hiện bất thường (AUC)}} &
\multirow{2}{*}{\textbf{MADRL Reward}} &
\multirow{2}{*}{\textbf{VAE MSE}} \\
\cmidrule(lr){3-4} \cmidrule(lr){5-7}
 & & \textbf{Accuracy} & \textbf{AUC} &
\textbf{LOF} & \textbf{Z-Score} & \textbf{DT-EKF} & & \\
\midrule
\multirow{3}{*}{Antwerp}
 & K-Means   & 0,7301 & 0,6755 & \multirow{3}{*}{0,506}
                                & \multirow{3}{*}{0,920}
                                & \multirow{3}{*}{\textbf{0,865}}
                                & \multirow{3}{*}{1916,01}
                                & \multirow{3}{*}{0,0082} \\
 & GMM       & 0,6558 & 0,6965 & & & & & \\
 & \textbf{DT-EKF} & \textbf{0,7403} & \textbf{0,7948} & & & & & \\
\midrule
\multirow{3}{*}{LoRaWAN}
 & K-Means   & 0,5050 & 0,5050 & \multirow{3}{*}{0,500}
                                & \multirow{3}{*}{0,919}
                                & \multirow{3}{*}{\textbf{0,834}}
                                & \multirow{3}{*}{2519,26}
                                & \multirow{3}{*}{0,0160} \\
 & GMM       & 0,5050 & 0,5164 & & & & & \\
 & \textbf{DT-EKF} & \textbf{0,7825} & \textbf{0,8614} & & & & & \\
\midrule
\multirow{3}{*}{Indoor}
 & K-Means   & 0,7673 & 0,8300 & \multirow{3}{*}{0,489}
                                & \multirow{3}{*}{0,958}
                                & \multirow{3}{*}{\textbf{0,978}}
                                & \multirow{3}{*}{3088,95}
                                & \multirow{3}{*}{0,0537} \\
 & GMM       & 0,6109 & 0,6657 & & & & & \\
 & \textbf{DT-EKF} & \textbf{0,9345} & \textbf{0,9725} & & & & & \\
\bottomrule
\end{tabular}%
}
\end{table}

% =============================================================================
\section{Thảo Luận}
% =============================================================================

\subsection{Ưu Điểm của DT-EKF so với Baseline}

\textbf{Phân loại chất lượng:} DT-EKF vượt trội K-Means và GMM về cả
Accuracy lẫn AUC trên toàn bộ ba bộ dữ liệu. Cải thiện nổi bật nhất là
trên LoRaWAN, nơi K-Means và GMM chỉ đạt AUC $\approx 0{,}50$ (tương đương
đoán ngẫu nhiên) do dữ liệu không có cấu trúc cụm rõ ràng trong không gian
đặc trưng tĩnh. DT-EKF đạt AUC = 0,861 và Accuracy = 0,783 nhờ khai thác
thông tin động học theo thời gian qua mô hình EKF. Trên Indoor, DT-EKF đạt
Accuracy = 0,935 và AUC = 0,973 — vượt K-Means lần lượt $+16{,}7\%$ và
$+14{,}3\%$.

\textbf{Phát hiện bất thường:} LOF hoàn toàn thất bại (AUC $\approx 0{,}49$--$0{,}51$)
do chỉ phân tích mật độ điểm dữ liệu trong không gian đặc trưng mà không
xét đến ngữ cảnh thời gian. Rolling Z-Score đạt hiệu quả tốt (AUC
$\approx 0{,}92$--$0{,}96$) nhờ tính chất tạm thời của lỗi tiêm, và nhỉnh
hơn DT-EKF trên Antwerp và LoRaWAN. Tuy nhiên, Z-Score chỉ xử lý một chiều
RSSI và không thể thực hiện phân loại chất lượng. Trong khi đó, DT-EKF
xử lý đồng thời cả hai chiều (RSSI và số trạm gốc) và thực hiện song song
cả hai nhiệm vụ, đạt AUC $\geq 0{,}978$ trên Indoor.

\subsection{Phân Tích Điểm Yếu và Giới Hạn}

\begin{itemize}
    \item \textbf{DT-EKF thua Z-Score về anomaly AUC trên Antwerp và LoRaWAN}:
    Rolling Z-Score cao hơn DT-EKF khoảng $5{,}5\%$ và $8{,}5\%$ tương ứng.
    Điều này phản ánh ưu thế của Z-Score khi phát hiện bất thường đơn biến
    (chỉ RSSI) trong khi EKF phải dung hòa cả hai chiều trạng thái trong
    ma trận nhiễu quan sát. Cải thiện đáng kể trên Indoor (DT-EKF 0,978 vs
    Z-Score 0,958) cho thấy ưu thế khi cả hai kênh đồng thời mang thông tin.

    \item \textbf{RSSI-RMSE LoRaWAN tương đối cao} (7,045 dBm): Tín hiệu
    LoRaWAN dao động trên dải rộng hơn, khiến xấp xỉ tuyến tính hóa cục bộ
    của EKF tại mỗi bước ít chính xác hơn so với Sigfox và WiFi.

    \item \textbf{Tính đại diện của fault injection}: Biên độ tiêm $\pm 2\sigma$
    phản ánh các lỗi thực tế ở mức vừa phải. Trong triển khai thực tế, các
    lỗi tinh tế hơn ($< 1\sigma$) hoặc lỗi drift dài hạn có thể đòi hỏi
    cơ chế phát hiện bổ sung.
\end{itemize}

% =============================================================================
\section{Kết Luận}
% =============================================================================

Thực nghiệm trên ba bộ dữ liệu IoT thực tế (Sigfox Antwerp, LoRaWAN Italy,
Indoor WiFi) chứng minh hiệu quả của khung đề xuất VAE + Digital Twin (EKF)
+ MADRL+GAT:

\begin{enumerate}
    \item \textbf{VAE} học biểu diễn tiềm ẩn có cấu trúc với MSE tái tạo
    thấp ($< 0{,}054$) trên cả ba bộ dữ liệu, cung cấp đầu vào chất lượng
    cho module MADRL.

    \item \textbf{DT-EKF} vượt trội toàn diện các baseline tĩnh (K-Means,
    GMM) về phân loại chất lượng trên mọi bộ dữ liệu. Cải thiện lớn nhất
    trên LoRaWAN: Accuracy tăng $+27{,}8\%$ và AUC tăng $+35{,}6\%$ so với
    K-Means. Về phát hiện bất thường, DT-EKF đạt AUC $\geq 0{,}834$ nhất
    quán và vượt Z-Score trên Indoor, trong khi cung cấp đồng thời khả năng
    phân loại chất lượng mà Z-Score không có.

    \item \textbf{MADRL+GAT} hội tụ ổn định và đơn điệu trong 50 episode
    trên cả ba bộ dữ liệu, với phần thưởng đánh giá tăng theo đặc tính
    môi trường IoT.
\end{enumerate}

\end{document}
